\section{Por qué se necesita Multi-Head Attention: un Weighted Average no basta}

La sección anterior presentó la fórmula; esta sección explica multi-head con la intuición del mecanismo que da el ponente. La attention de una sola cabeza produce, al final, un solo promedio ponderado; si todas las relations comparten la misma value projection, las relaciones sintácticas, posicionales, de entidades o de copia compiten por el mismo espacio de representación. Multi-head usa varios grupos de projections para aprender en paralelo distintos enrutamientos y transformaciones.

\subsection{La Attention de una sola cabeza es un promedio ponderado con una Value Transform compartida}

Para la posición $i$, la salida de una sola cabeza puede escribirse así:
\[
y_i=\sum_{j=1}^{n}\alpha_{ij}W_Vx_j,
\qquad \sum_j\alpha_{ij}=1,\quad \alpha_{ij}\ge 0.
\]
Aquí $W_V$ es compartida por todas las posiciones que se leen, y $\alpha_{ij}$ cambia según el contenido de query/key. El enrutamiento es dinámico, pero el contenido transmitido pasa por la misma transformación lineal; por eso la relation capacity de un solo head es limitada.

\lecturefigure{slide-21-attention-weighted-average.jpg}{La attention de una sola cabeza puede verse como un promedio ponderado dinámico de las values}{Diapositiva V021 recuperada de la grabación oficial; 00:25:30}

\begin{knowledgebox}{Lectura de la figura: lo dinámico son los pesos; lo fijo es la transformación dentro de un mismo head}
El grosor de cada flecha lo determina $\alpha_{ij}$, y todas las entradas pasan primero por la misma $W_V$. Observe primero que puede tomar información de cualquier posición y después si los distintos offsets o relaciones tienen una transform independiente; esto explica por qué una sola cabeza es más flexible que un pooling fijo, pero aun así puede comprimir varios tipos de relación en el mismo subespacio.
\end{knowledgebox}

\subsection{La Convolution ofrece una Transform distinta para cada Offset}

La salida de una convolution unidimensional es aproximadamente:
\[
y_i=\sum_{\delta\in\mathcal{K}}W_{\delta}x_{i+\delta}.
\]
Aquí $\mathcal{K}$ son los kernel offsets, y $W_{\delta}$ puede variar con la posición relativa. En la convolution, las posiciones de conexión son fijas y las transforms son diversas; en la single-head attention, los pesos de conexión son dinámicos y la value transform es compartida. Multi-head busca obtener a la vez enrutamiento dinámico y varios grupos de transforms.

\lecturefigure{slide-22-convolution-linear-transformations.jpg}{La convolution usa una linear transformation distinta para cada posición relativa}{Diapositiva V022 recuperada de la grabación oficial; 00:27:03}

\begin{knowledgebox}{Lectura de la figura: dónde coloca sus grados de libertad la Attention y dónde la Convolution}
Las aristas de colores representan el kernel weight de cada offset. La convolución puede distinguir «una posición a la izquierda» de «dos posiciones a la derecha», pero el alcance de lectura lo predefine el kernel; la attention puede elegir cualquier posición según el contenido, pero dentro de una sola cabeza comparte la value transform. Solo al entender este contraste se ve que multi-head no es un simple ensemble.
\end{knowledgebox}

\subsection{Multi-Head: varios Relation Subspaces trabajan en paralelo}

La \term{multi-head attention (atención multicabeza)} usa para cada head sus propias $W_Q^h,W_K^h,W_V^h$ y después concatena las salidas:
\[
\operatorname{MHA}(X)=\operatorname{Concat}(H_1,\ldots,H_H)W_O,
\qquad H_h=\operatorname{Attn}(XW_Q^h,XW_K^h,XW_V^h).
\]
Aquí $H$ es el head count y $W_O$ es la projection de salida. Cada head puede aprender relaciones locales, copia, correferencia, posición u otras interactions, pero no hay garantía de que cada head corresponda a un concepto que una persona pueda nombrar.

\lecturefigure{slide-23-multi-head-attention.jpg}{La multi-head attention ofrece varios grupos de enrutamiento y de value transformation}{Diapositiva V023 recuperada de la grabación oficial; 00:28:15}

\begin{importantbox}{Lo esencial de multi-head no es «mirar varias veces»}
Cada head cambia a la vez el query/key routing y el value feature space. Esto permite al modelo expresar varios tipos de relación en la misma capa y escribir el resultado de vuelta en el residual stream compartido; más heads no es automáticamente mejor: la head dimension, la redundancia, la KV memory y el serving cost determinan en conjunto la escala efectiva.
\end{importantbox}

\begin{lstlisting}[caption={Estructura conceptual de la multi-head attention y de la KV compartida}]
for head in query_heads:
    query = project_q(hidden, head)
    kv_group = head_to_kv_group(head)  # MHA: one group/head; GQA/MQA: shared
    key = project_k(hidden, kv_group)
    value = project_v(hidden, kv_group)
    head_output[head] = attention(query, key, value, causal=True)
output = project_o(concat(head_output))
\end{lstlisting}

\teachervoice{El ponente resume la single-head attention como «un solo weighted average» y luego usa las offset-specific transforms de la convolution para explicar por qué se necesitan varias cabezas. Esta explicación del mecanismo es más precisa que decir que «varias cabezas pueden atender a posiciones distintas»: las cabezas múltiples aportan a la vez distintos espacios de similitud y distintas transformaciones del contenido.}

\subsection{Resumen de la sección}

La single-head attention combina enrutamiento dinámico con una value transform compartida, mientras que la convolution ofrece enrutamiento fijo con offset-specific transforms; la multi-head attention ejecuta en paralelo varios grupos de enrutamiento dinámico y de feature transform. Más adelante, MQA/GQA volverán a decidir qué projections deben ser independientes para reducir la KV memory.
